\documentclass{situationist-preprint}

\title{Deadlines and Self-Control}
\msid{ACT-13-b}
\msfamily{ACT-13}
\msdiscipline{Action, Agency and Free Will}
\mstype{Companion manuscript}
\mskeywords{procrastination; deadlines; self-control; precommitment; hyperbolic discounting; personal rules}

\begin{document}
\maketitle

\begin{abstract}
ACT-13-a argues that procrastination is a form of akrasia involving temporal discounting. We draw the practical consequence for one familiar remedy. A deadline is an instrument of self-control: by attaching a proximate cost to further delay, it brings the motivational pull of acting into line with the agent's better judgment at the moments at which discounting would otherwise defeat it. We show why deadlines work on the account of ACT-13-a and why they would be pointless on rival accounts, why self-imposed deadlines are weaker than external ones, and why the agent who adopts a deadline against her own foreseen procrastination acts rationally, although her need for it is a mark of irrationality.
\end{abstract}

\section{The problem}

People who procrastinate set themselves deadlines, ask others to set them, and work better when they have them. \citet{Ariely2002} found that students permitted to set their own deadlines for course papers chose binding deadlines before the end of term, at a cost to themselves, and performed better than students with a single final deadline, though worse than students given evenly spaced external deadlines. Economists have modeled such behavior as the commitment of a sophisticated agent who foresees that her future preferences will diverge from her present ones \citep{Strotz1955,ODonoghueRabin1999}, and \citet{Elster1979} made Ulysses binding himself to the mast the paradigm of rational precommitment.

The account of procrastination defended in ACT-13-a explains why deadlines work, and its explanation differs from the economists' in one respect that matters. On the economic models the agent's preferences reverse; on ACT-13-a her motivations reverse while her judgment does not. Procrastination is strict akrasia (Lemma~3.1 of ACT-13-a): at the moment of delay, the procrastinator judges that beginning now is better, all things considered. A deadline is therefore not a device by which one self constrains a later self with different values, but a device by which a single agent with stable values protects her conduct from a foreseen failure of motivation.

\section{How deadlines work}

\begin{definition}[Deadline]\label{d:deadline}
A \emph{deadline} for a task is a time $d$ together with a cost $c$ that the agent will incur if the task is not completed by $d$. It is \emph{external} if the cost is imposed and enforced by others, and \emph{self-imposed} if it is imposed by the agent and its enforcement depends on her.
\end{definition}

\begin{proposition}\label{p:works}
A deadline at $d$ with cost $c$ removes the occasion for procrastination at moments close enough to $d$ that the discounted cost of missing it exceeds the discounted cost of beginning.
\end{proposition}

\begin{proof}
By Lemma~3.2 of ACT-13-a, the procrastinator delays because the motive to delay, which is supplied by the proximate costs of the task, exceeds the motive supplied by her judgment, whose object, the benefit of completion, is remote and so heavily discounted (Definition~2.2 of ACT-13-a). A deadline adds to the motive to begin a further consideration, the avoidance of $c$, whose time is $d$. As $d$ approaches, the discount on $c$ shrinks steeply, since the discounting is hyperbolic. At some moment the discounted cost of missing the deadline exceeds the discounted cost of beginning, and from then on the motivationally strongest option is to begin. Since beginning is also what she judges best, the gap between judgment and motive that constitutes her akrasia is closed for those moments.
\end{proof}

Two features of the proof deserve note. First, the deadline does not alter the agent's judgment. By Principle~P1 of ACT-13-a her judgment of the value of beginning does not vary with proximity, and the deadline adds nothing to what she already judged; what it changes is the temporal location of a cost, and so its motivational weight. Second, a single deadline helps only near $d$. Between the adoption of the deadline and the moment at which its cost becomes motivationally decisive, the structure of ACT-13-a is unchanged, and the agent may procrastinate as before. This is the familiar last-minute rush, and it explains the finding of \citet{Ariely2002} that evenly spaced deadlines outperform a single one: they bring a proximate cost to bear at more of the moments at which she could begin.

\begin{proposition}\label{p:pointless}
On an account on which procrastination is the rational expression of a preference for the present, a deadline adopted to prevent procrastination is irrational.
\end{proposition}

\begin{proof}
Suppose that the procrastinator's delay at each moment expresses her considered preference at that moment, as some theories of rational choice permit. Then the deadline is adopted at an earlier moment to defeat her considered preference at later moments, at a cost $c$ risked for nothing she will then value. If each moment's preference is authoritative for that moment, the earlier self has no standing to overrule the later ones, and imposing a cost to do so is a loss by the standards of the later selves on whom it falls.
\end{proof}

The rival account must therefore treat the practice of setting deadlines, which procrastinators adopt freely and later endorse, as a mistake. The account of ACT-13-a treats it as the rational response of an agent who foresees that her motivation will fail to track her judgment, which is what it appears to be.

\section{Self-imposed deadlines}

A self-imposed deadline differs from an external one in that its cost depends on the agent's own enforcement. If, when the deadline comes, she can waive it at no cost, it adds nothing to the motive to begin, and Proposition~\ref{p:works} does not apply. Self-imposed deadlines work, when they do, because the agent finds a way to attach a cost to waiving them.

\begin{proposition}\label{p:self}
A self-imposed deadline is effective only to the extent that waiving it carries a cost that is not discounted away at the moment of waiving.
\end{proposition}

\begin{proof}
At $d$ the agent who has not completed the task faces a choice between incurring $c$ and waiving the deadline. If waiving is costless, she waives it, and the deadline exerts no motivational force at earlier moments either, since she foresees the waiver. The deadline can exert force only if waiving costs her something whose discounted weight at $d$ is sufficient to deter her.
\end{proof}

There are two sources of such costs. One is social: the agent announces the deadline, deposits a sum she will forfeit, or asks another to hold her to it, and so turns a self-imposed deadline into a partly external one. The other is internal. \citet{Ainslie1992} argues that an agent who treats each choice as a precedent for a whole series of similar choices, under a personal rule, weighs the sum of the series' discounted values rather than the value of the single choice, and that the summed series favors the larger, later rewards. Waiving a deadline then costs her the credibility of the rule, which is a cost now, not later. The effectiveness of self-imposed deadlines thus depends on resolve, and \citet{Ariely2002}'s finding that they are weaker than external ones is what Proposition~\ref{p:self} predicts.

\section{Objections and replies}

\begin{objection}
An agent who needs a deadline to do what she judges best is not exercising self-control but outsourcing it; the self-controlled agent acts on her judgment without external help.
\end{objection}

\begin{reply}
Self-control is the agent's capacity to bring her conduct into line with her better judgment against contrary motivation \citep{Elster1979}. It may be exercised synchronically, by effort at the moment of temptation, or diachronically, by arranging in advance that the temptation will be weaker or will not arise. The agent who sets a deadline exercises the diachronic form: she acts at the moment of adoption on her judgment about her later conduct, and the deadline is her instrument. That she needs it is, by Theorem~3.3 of ACT-13-a, a mark of a disposition to irrationality; that she adopts it is rational.
\end{reply}

\begin{objection}
Deadlines produce work done hastily at the last moment, which is often worse than work done without a deadline.
\end{objection}

\begin{reply}
That is the consequence of a single deadline noted after Proposition~\ref{p:works}, and it is a reason to space deadlines, not to abandon them. Where the agent would not procrastinate, no deadline is needed and none may help; the claim is only that, given the discounting that ACT-13-a identifies, a deadline closes the gap between judgment and motive at the moments it reaches.
\end{reply}

\section{Conclusion}

If procrastination is akrasia driven by temporal discounting, as ACT-13-a argues, deadlines are instruments of self-control. They do not change what the agent judges best; they move a cost close enough in time to outweigh the discounted costs of beginning. Spaced deadlines work better than single ones because they reach more of the moments at which the agent could begin, and self-imposed deadlines work only when waiving them is costly. The agent who binds herself acts rationally, in response to a foreseen irrationality.

\begin{thebibliography}{9}\small
\bibitem[Ainslie(1992)]{Ainslie1992} Ainslie, G. (1992). \emph{Picoeconomics: The Strategic Interaction of Successive Motivational States within the Person}. Cambridge University Press.
\bibitem[Ariely and Wertenbroch(2002)]{Ariely2002} Ariely, D. and Wertenbroch, K. (2002). Procrastination, deadlines, and performance: self-control by precommitment. \emph{Psychological Science}, 13(3), 219--224.
\bibitem[Elster(1979)]{Elster1979} Elster, J. (1979). \emph{Ulysses and the Sirens: Studies in Rationality and Irrationality}. Cambridge University Press.
\bibitem[O'Donoghue and Rabin(1999)]{ODonoghueRabin1999} O'Donoghue, T. and Rabin, M. (1999). Doing it now or later. \emph{American Economic Review}, 89(1), 103--124.
\bibitem[Strotz(1955)]{Strotz1955} Strotz, R.~H. (1955). Myopia and inconsistency in dynamic utility maximization. \emph{The Review of Economic Studies}, 23(3), 165--180.
\end{thebibliography}

\end{document}
